% =========================================================
% Proof file: prf-unique-additive-inverse.tex
% =========================================================

\newpage
\phantomsection
\label{prf:unique-additive-inverse}

\begin{remark*}[Return]
\hyperref[thm:unique-additive-inverse]{Return to Theorem}
\end{remark*}

\begin{theorem*}[Unique Additive Inverse]
For every $v \in V$, if $v_0,v_1 \in V$ satisfy
\[
v+v_0=0
\qquad\text{and}\qquad
v+v_1=0,
\]
then
\[
v_0=v_1.
\]
\end{theorem*}

\begin{proof}
\textbf{Professional Standard Proof.}~\\
Let $v \in V$, and let $v_0,v_1 \in V$ satisfy
\[
v+v_0=0
\qquad\text{and}\qquad
v+v_1=0.
\]
Add $v_1$ to both sides of $v+v_0=0$:
\[
v_1+(v+v_0)=v_1+0.
\]
By associativity,
\[
(v_1+v)+v_0=v_1.
\]
By commutativity,
\[
v_1+v=v+v_1=0.
\]
Hence
\[
0+v_0=v_1,
\]
so
\[
v_0=v_1.
\]
Therefore the additive inverse of $v$ is unique.
\end{proof}

\begin{proof}
\textbf{Detailed Learning Proof.}~\\

\textbf{Step 1.} Fix a vector and two candidate inverses.\\
Let $v \in V$, and suppose $v_0,v_1 \in V$ satisfy
\[
v+v_0=0
\qquad\text{and}\qquad
v+v_1=0.
\]
The goal is to show that these two candidates are equal.

\textbf{Step 2.} Start from one inverse equation and add the other inverse to both sides.\\
From
\[
v+v_0=0,
\]
add $v_1$ to both sides:
\[
v_1+(v+v_0)=v_1+0.
\]

\textbf{Step 3.} Simplify the right-hand side.\\
Because $0$ is the additive identity,
\[
v_1+0=v_1.
\]
So
\[
v_1+(v+v_0)=v_1.
\]

\textbf{Step 4.} Reassociate the left-hand side.\\
By associativity of addition,
\[
(v_1+v)+v_0=v_1.
\]

\textbf{Step 5.} Use commutativity to bring the known inverse relation into view.\\
By commutativity,
\[
v_1+v=v+v_1.
\]
But $v_1$ is an additive inverse of $v$, so
\[
v+v_1=0.
\]
Hence
\[
(v_1+v)+v_0=0+v_0.
\]
Therefore
\[
0+v_0=v_1.
\]

\textbf{Step 6.} Finish with the identity law.\\
Because $0$ is the additive identity,
\[
0+v_0=v_0.
\]
So
\[
v_0=v_1.
\]

\textbf{Step 7.} State the conclusion.\\
Any two additive inverses of the same vector must agree. Therefore every vector has a unique additive inverse.
\end{proof}

\begin{remark*}[Proof structure]
The proof is direct. One inverse equation is modified by adding the other inverse, and the vector space axioms reduce the resulting expression until the two candidates are shown to be equal.
\end{remark*}

\begin{remark*}[Dependencies]
\hfill
\begin{itemize}
  \item \hyperref[def:vector-space]{Definition of Vector Space}
  \item \hyperref[thm:unique-additive-identity]{Unique Additive Identity}
\end{itemize}
\end{remark*}